% ============================================================
%  ANNEXE UNIQUE : EXTRAITS DE CODE SOURCE
% ============================================================
\appendix

\chapter{Extraits de code source}
\label{ann:code}

Cette annexe contient les extraits les plus significatifs du code d'entraînement, de préparation des données et de guidage du modèle développé dans ce projet.

\subsection{Configuration et Hyperparamètres}

L'extrait suivant présente les hyperparamètres et la configuration globale utilisés pour initialiser l'entraînement (définis dans la cellule 21 du notebook).

\begin{lstlisting}[language=Python, caption={Hyperparamètres globaux d'entraînement.}]
# -- Image
IMG_SIZE    = 320
IN_CHANNELS = 3

# -- Training
EPOCHS         = 150
BATCH_SIZE     = 8
LEARNING_RATE  = 3e-4
WEIGHT_DECAY   = 1e-4
MAX_GRAD_NORM  = 1.0

# -- Scheduler & early stop
PATIENCE_ES    = 15
PATIENCE_LR    = 5

# -- ImageNet stats (correct for pretrained encoders)
IMAGENET_MEAN = (0.485, 0.456, 0.406)
IMAGENET_STD  = (0.229, 0.224, 0.225)
\end{lstlisting}

\subsection{Dataset et augmentation de données}
\label{sec:code-dataset-augmentation}

L'implémentation de la classe personnalisée PyTorch \texttt{DroneDataset} et la chaîne d'augmentation de données basée sur Albumentations (cellules 27 et 29 du notebook).

\begin{lstlisting}[language=Python, caption={Pipeline d'augmentation et classe de chargement des données.}]
# Pipeline d'augmentation pour l'entraînement
train_transform = A.Compose([
    A.Resize(IMG_SIZE, IMG_SIZE),
    A.HorizontalFlip(p=0.5),
    A.VerticalFlip(p=0.3),
    A.RandomRotate90(p=0.5),
    A.ShiftScaleRotate(
        shift_limit=0.1, scale_limit=0.2, rotate_limit=30,
        border_mode=cv2.BORDER_REFLECT_101, p=0.7
    ),
    A.Perspective(scale=(0.04, 0.08), p=0.3),
    A.RandomBrightnessContrast(brightness_limit=0.3, contrast_limit=0.3, p=0.6),
    A.CoarseDropout(
        max_holes=8, max_height=32, max_width=32,
        min_holes=1, fill_value=0, mask_fill_value=0, p=0.3
    ),
    A.Normalize(mean=IMAGENET_MEAN, std=IMAGENET_STD),
    ToTensorV2(),
])

class DroneDataset(Dataset):
    """Binary segmentation: SAFE (1) vs UNSAFE (0) landing zones."""
    SAFE_CLASSES = [1, 3] # paved-area (1), grass (3)

    def __init__(self, img_list, mask_list, img_dir, mask_dir, transform=None):
        self.img_list  = img_list
        self.mask_list = mask_list
        self.img_dir   = img_dir
        self.mask_dir  = mask_dir
        self.transform = transform

    def __len__(self):
        return len(self.img_list)

    def __getitem__(self, idx):
        img_path  = os.path.join(self.img_dir,  self.img_list[idx])
        mask_path = os.path.join(self.mask_dir, self.mask_list[idx])

        image = cv2.imread(img_path)
        image = cv2.cvtColor(image, cv2.COLOR_BGR2RGB)
        mask = cv2.imread(mask_path, cv2.IMREAD_GRAYSCALE)

        # Binarisation: 1.0 = SAFE, 0.0 = UNSAFE
        mask = np.isin(mask, self.SAFE_CLASSES).astype(np.float32)

        if self.transform:
            augmented = self.transform(image=image, mask=mask)
            image = augmented["image"]
            mask  = augmented["mask"]

        return image, mask.unsqueeze(0)
\end{lstlisting}

\subsection{Architecture du modèle et fonction de coût}

Le code ci-dessous montre la définition du modèle via la bibliothèque \texttt{segmentation\_models\_pytorch} et la fonction de coût hybride combinant la BCE logarithmique et la Tversky Loss.

\begin{lstlisting}[language=Python, caption={Définition du modèle U-Net et de la fonction de coût.}]
# Modèle U-Net avec encodeur EfficientNet-B0 pré-entraîné
model = smp.Unet(
    encoder_name    = "efficientnet-b0",
    encoder_weights = "imagenet",
    in_channels     = IN_CHANNELS,
    classes         = 1,
    activation      = None,
    decoder_use_batchnorm = True,
)
model = model.to(DEVICE)

# Definition des fonctions de cout
bce_loss = nn.BCEWithLogitsLoss()
tversky_loss = smp.losses.TverskyLoss(
    mode="binary", alpha=0.3, beta=0.7, smooth=1e-6
)

def loss_fn(pred, target):
    return 0.5 * bce_loss(pred, target) + 0.5 * tversky_loss(pred, target)
\end{lstlisting}

\subsection{Contrôleur de guidage et machine d'états}
\label{sec:code-guidage-machine}

Ces classes et fonctions implémentent le pipeline d'évitement géométrique et la logique de pilotage de haut niveau qui traite la prédiction du modèle.

\begin{lstlisting}[language=Python, caption={Algorithme de guidage géométrique et machine d'états de vol.}]
# Guidance Controller
def guidance_controller(landing_x, landing_y, image_width, image_height, tolerance=10):
    center_x = image_width // 2
    center_y = image_height // 2

    # Calcul de l'erreur en pixels
    dx = landing_x - center_x
    dy = landing_y - center_y

    distance = math.sqrt(dx**2 + dy**2)
    heading = math.degrees(math.atan2(dy, dx))

    # Logique de guidage horizontal et vertical
    horizontal = "MOVE RIGHT" if dx > tolerance else ("MOVE LEFT" if dx < -tolerance else "HOLD")
    vertical = "MOVE BACKWARD" if dy > tolerance else ("MOVE FORWARD" if dy < -tolerance else "HOLD")
    centered = abs(dx) <= tolerance and abs(dy) <= tolerance

    return {
        "landing_x": landing_x, "landing_y": landing_y,
        "center_x": center_x, "center_y": center_y,
        "dx": dx, "dy": dy, "distance": distance, "heading": heading,
        "horizontal": horizontal, "vertical": vertical, "centered": centered
    }

# Machine d'états de vol
class FlightState(Enum):
    SEARCHING = "SEARCHING"
    GUIDING = "GUIDING"
    LANDING = "LANDING"

class FlightStateMachine:
    def __init__(self):
        self.state = FlightState.SEARCHING

    def update(self, landing_allowed, centered):
        if self.state == FlightState.SEARCHING:
            if landing_allowed:
                self.state = FlightState.GUIDING
        elif self.state == FlightState.GUIDING:
            if not landing_allowed:
                self.state = FlightState.SEARCHING
            elif centered:
                self.state = FlightState.LANDING
        elif self.state == FlightState.LANDING:
            if not landing_allowed:
                self.state = FlightState.SEARCHING
        return self.state
\end{lstlisting}

\clearpage
